\chapter{Introducci\'on}

El an\'alisis de im\'agenes m\'edicas es un campo donde la ingenier\'ia inform\'atica puede aportar herramientas de apoyo para tareas de segmentaci\'on, medici\'on, visualizaci\'on y documentaci\'on. En particular, la resonancia magn\'etica (RM) de columna lumbar permite observar estructuras como cuerpos vertebrales, discos intervertebrales y canal espinal, cuyo an\'alisis requiere criterio profesional y una correcta interpretaci\'on anat\'omica. En este contexto existe una oportunidad de asistencia inform\'atica: automatizar tareas operativas que hoy pueden requerir revisi\'on manual o semimanual, sin sustituir la decisi\'on del especialista. La segmentaci\'on anat\'omica asistida por inteligencia artificial puede servir como base para visualizar regiones de inter\'es, calcular mediciones geom\'etricas simples y organizar resultados de manera trazable.

El presente Proyecto Final de Ingenier\'ia propone el dise\~no y desarrollo de un prototipo funcional de software para RM lumbar sagital. El MVP contempla la carga y normalizaci\'on de estudios, la segmentaci\'on de v\'ertebras, discos intervertebrales y canal espinal, la extracci\'on de mediciones cuantitativas acotadas, la visualizaci\'on superpuesta de m\'ascaras y la generaci\'on de una salida estructurada editable. El sistema se plantea como herramienta de apoyo: no emite diagn\'osticos, no recomienda tratamientos y no reemplaza al profesional. El desarrollo se apoya principalmente en datasets p\'ublicos anotados, especialmente SPIDER, lo que favorece la reproducibilidad acad\'emica y permite validar t\'ecnicamente las segmentaciones contra m\'ascaras de referencia. Esta decisi\'on tambi\'en reduce la dependencia inicial de datos cl\'inicos privados y mantiene el alcance dentro de una prueba de concepto verificable.

La organizaci\'on del documento avanza desde la definici\'on del problema hacia la fundamentaci\'on y el dise\~no del prototipo. Los primeros cap\'itulos presentan el problema, la justificaci\'on, los objetivos, el alcance y la soluci\'on propuesta. Luego se desarrollan el estado del arte y el marco te\'orico, que establecen los antecedentes y conceptos necesarios para comprender las decisiones t\'ecnicas del proyecto. En etapas posteriores del documento final se incorporar\'an user research, requerimientos, arquitectura, metodolog\'ia, implementaci\'on, validaci\'on, resultados y conclusiones.
